\documentclass[aspectratio=169,11pt]{beamer}

\usetheme{default}
\usecolortheme{seahorse}
\setbeamertemplate{navigation symbols}{}
\setbeamertemplate{footline}[frame number]
\setbeamertemplate{itemize items}[circle]
\setbeamercolor{title}{fg=black}
\setbeamercolor{frametitle}{fg=black}
\setbeamerfont{frametitle}{series=\bfseries}

\usepackage{amsmath,amssymb}
\usepackage{booktabs}

\newcommand{\R}{\mathbb{R}}
\newcommand{\GP}{\mathcal{GP}}
\newcommand{\Norm}{\mathcal{N}}
\newcommand{\Ulost}{U_{\text{lost}}}
\newcommand{\good}{\textcolor{green!50!black}}
\newcommand{\bad}{\textcolor{red!70!black}}

\title{\bfseries Knowing the Outer Function vs.\ Observing the Intermediates}
\subtitle{Splitting the composite-BO gain into its two parts}
\date{}
\author{}

\begin{document}

\begin{frame}[plain]
\titlepage
\vfill
\small
\texttt{github.com/MuhammadAli-CS/three-arm-composite}
\end{frame}

% ------------------------------------------------------------------
\begin{frame}{Where we are}

Composite problems: $f(x)=g(h(x))$, with $h$ an expensive black box and $g$ a
known, cheap function.

\medskip
\begin{itemize}
\item \textbf{Direct}: fit a GP to $f$, ignore the structure.
\item \textbf{Composite}: fit GPs to $h$, push samples through $g$.
\end{itemize}

\medskip
Composite won most benchmarks --- by anywhere from $\mathbf{121\%}$ to
\textbf{nothing}.

\bigskip
\begin{block}{The question this work answers}
We do not know \emph{why} it wins, because composite changes two things at once.
\end{block}

\end{frame}

% ------------------------------------------------------------------
\begin{frame}{The confound}

Going from direct to composite gives the surrogate \textbf{two} advantages
simultaneously:

\bigskip
\begin{enumerate}
\item \textbf{It knows $g$} --- handed the outer function, never learns it from data.
\item \textbf{It observes $h$} --- each expensive run returns $k$ numbers, not one.
\end{enumerate}

\bigskip
Every result we have bundles these together.

\begin{itemize}
\item Where composite wins big --- which advantage did it?
\item Where composite wins nothing --- which one failed?
\end{itemize}

\bigskip
\centering
\emph{Two numbers cannot separate two effects.}

\end{frame}

% ------------------------------------------------------------------
\begin{frame}{The fix: a third arm}

\begin{center}
\begin{tabular}{lcc}
\toprule
 & $g$ unknown & $g$ known \\
\midrule
observes $y$ only & arm 1 \textbf{direct} & \alert{arm 3 \textbf{latent}} \\
observes $h$      & --- (incoherent)      & arm 2 \textbf{composite} \\
\bottomrule
\end{tabular}
\end{center}

\bigskip
\textbf{Arm 3} knows $g$ exactly, holds the \emph{same} GP prior on $h$ as
composite, but never sees $h$ --- it must infer it from $y$.

\bigskip
\begin{align*}
\text{latent} - \text{direct}    &= \text{the value of \textbf{knowing} } g\\
\text{composite} - \text{latent} &= \text{the value of \textbf{observing} } h
\end{align*}

\medskip
The two sum to the composite gain we already publish --- so the decomposition
must add up, and any failure is a bug, not a finding.

\end{frame}

% ------------------------------------------------------------------
\begin{frame}{What should govern the split}

Observing $h$ can only help where $y=g(h)$ fails to pin $h$ down.
\textbf{Invertibility of $g$ is the controlling property.}

\bigskip
\begin{columns}[T]
\begin{column}{0.48\textwidth}
\textbf{$g$ injective}\\[2pt]
$h=g^{-1}(y)$, so arms 2 and 3 see identical information.\\[4pt]
$\Rightarrow$ observing $h$ worth \textbf{exactly 0}
\end{column}
\begin{column}{0.48\textwidth}
\textbf{$g$ linear}\\[2pt]
A sum of independent GPs \emph{is} a GP, so arm 3 collapses to arm 1.\\[4pt]
$\Rightarrow$ knowing $g$ worth \textbf{$\approx 0$}
\end{column}
\end{columns}

\bigskip
\begin{block}{Why the linear case matters for the paper}
Truss and RCM40 both have linear outer maps --- and both gave our weakest
results.
\end{block}

\bigskip
In between: what $g$ destroys is the \emph{fiber} $g^{-1}(y)$ --- two points for
$z^2$, several for $\cos$, a half-line for a clipped map, a circle for a norm.

\end{frame}

% ------------------------------------------------------------------
\begin{frame}{How arm 3 infers $h$ --- the whole idea in one example}

$g(z)=z^2$. You observe $y=4$ at \textbf{two nearby inputs}.

\medskip
Four candidates fit the data \emph{perfectly}:
$(+2,+2)$, $(-2,-2)$, $(+2,-2)$, $(-2,+2)$.

\medskip
The data cannot choose. \textbf{Smoothness can}: a jump from $+2$ to $-2$ over a
tiny step is not a plausible function.

\bigskip
\begin{center}
\begin{tabular}{lrrr}
\toprule
distance between inputs & correlation & same-sign advantage \\
\midrule
close   & $0.9$  & $2.9\times10^{16}$ \\
medium  & $0.3$  & $14\times$ \\
far     & $0.02$ & $1.2\times$ \\
\bottomrule
\end{tabular}
\end{center}

\bigskip
Note what is learned: \emph{not} whether $h=+2$ or $-2$ (still 50/50), but that
both points did the \textbf{same} thing.

\end{frame}

% ------------------------------------------------------------------
\begin{frame}{That is Bayes' rule}

Score every candidate filling-in of $h$ on two questions:

\begin{equation*}
p(\mathbf h \mid \mathbf y)\;\propto\;
\underbrace{\prod_{i=1}^{n}\Norm\bigl(y_i;\,g(h(x_i)),\,\sigma^2\bigr)}_{\text{does it match the data?}}
\;\cdot\;
\underbrace{\prod_{j=1}^{k}\Norm\bigl(h_j(X);\,m_j,\,v_jK\bigr)}_{\text{is it a plausible function?}}
\end{equation*}

\bigskip
Neither factor works alone: the likelihood cannot choose between $\pm2$; the
prior does not know the data.

\bigskip
\begin{alertblock}{Why this is the hard part}
Arms 1 and 2 are ordinary GP regression --- closed form.
A Gaussian prior times a \emph{non}-Gaussian likelihood is not Gaussian, so arm 3
has \textbf{no formula} and needs MCMC. Every difficulty in this project has
been here.
\end{alertblock}

\end{frame}

% ------------------------------------------------------------------
\begin{frame}{What the two factors actually compute}

First, the notation. $\Norm(a;\,b,\,c)$ is not a distribution --- it is a
\textbf{number}: the height of the bell curve with mean $b$ and variance $c$,
read off at the position $a$. Three numbers in, one number out.

\bigskip
\begin{columns}[T]
\begin{column}{0.50\textwidth}
\textbf{Factor 1 --- does it match?}\\[4pt]
Centre the bell on what the candidate \emph{predicts}, $g(h(x_i))$; read it off
at what you \emph{measured}, $y_i$.\\[5pt]
$\sigma$ is the tolerance, and it sits in the denominator of the exponent --- at
$\sigma=0.01$, missing by $0.1$ costs $e^{-50}$.\\[5pt]
$\textstyle\prod_i$ runs over evaluated points: once $\mathbf h$ is fixed, the
measurements are independent.
\end{column}
\begin{column}{0.46\textwidth}
\textbf{Factor 2 --- is it smooth?}\\[4pt]
One multivariate Gaussian over the \emph{whole vector} of values, with
$K_{ii'}=\exp(-\|x_i-x_{i'}\|^2/2\ell^2)$.\\[3pt]
Three points at $x=(0,\,0.1,\,0.6)$, $\ell=0.5$:
\[
K=\begin{pmatrix}1 & 0.98 & 0.49\\ 0.98 & 1 & 0.61\\ 0.49 & 0.61 & 1\end{pmatrix}
\]
$K^{-1}$ in that exponent \textbf{is} the smoothness machinery.
\end{column}
\end{columns}

\end{frame}

% ------------------------------------------------------------------
\begin{frame}{Evaluating it on three candidates}

$g(z)=z^2$, $\sigma=0.01$, measured $\mathbf y=(4,\;4.2025,\;1)$.

\medskip
\begin{center}
\small
\begin{tabular}{llrrr}
\toprule
Candidate & $\mathbf h$ & log-likelihood & log-prior & log-posterior \\
\midrule
\textbf{A} smooth, matches & $(2,\;2.05,\;1)$            & $+11.06$      & $-2.80$        & \good{$+8.26$} \\
\textbf{B} jumpy, matches  & $(2,\;\mathbf{-2.05},\;1)$  & $+11.06$      & \bad{$-469.5$} & $-458.45$ \\
\textbf{C} smooth, misses  & $(2,\;2.02,\;\mathbf{1.5})$ & \bad{$-7876$} & $-2.80$        & $-7878.8$ \\
\bottomrule
\end{tabular}
\end{center}

\bigskip
\textbf{A and B have the \emph{identical} log-likelihood} --- not similar,
identical, because $(-2.05)^2=(2.05)^2$. That factor is mathematically incapable
of separating them. The prior does it alone: $\mathbf h^\top K^{-1}\mathbf h$ is
$4.42$ for A and $937.8$ for B, because $K_{12}=0.98$ demands that the first two
points agree.

\medskip
C inverts it --- a marginally \emph{better} prior than A, destroyed by predicting
$2.25$ where you measured $1$.

\bigskip
\centering
\emph{Each candidate is killed by a different factor. That is why the equation
is a product.}

\end{frame}

% ------------------------------------------------------------------
\begin{frame}{The experiment}

19 outer maps applied to inner functions \textbf{drawn from the same GP family
the surrogates assume}, so $h$ is never a confound.

\medskip
\begin{center}
\small
\begin{tabular}{lll}
\toprule
Group & Maps & Purpose \\
\midrule
Controls & $z$, \; $e^{z}$, \; $\sinh z$, \; $y_1+y_2$ & answers known in advance \\
Sweep & $(z-c)^2$, $c=0,0.5,1,2,3$ & one knob on \emph{how much} is destroyed \\
Pairs & $|z|$, $\cos(\omega z)$, censoring maps & one property changed at a time \\
Analogs & sum, product, norm, max, ratio & mirror Truss, DTLZ2, Color, Lake, RGB \\
\bottomrule
\end{tabular}
\end{center}

\bigskip
$d=6$, 5 initial points $+$ 40 evaluations, \textbf{4 function draws $\times$ 5
designs}, all arms sharing the initial design and acquisition.

\medskip
$\Rightarrow$ 380 BO runs per arm, $\approx$ 40 min on 16 cores.

\end{frame}

% ------------------------------------------------------------------
\begin{frame}{Controls are not optional}

The first full grid looked \textbf{completely reasonable} --- smooth trends,
sensible signs, tight error bars.

\medskip
It was wrong.

\bigskip
\begin{center}
\texttt{identity}: ``observing $h$'' must be \textbf{exactly 0}.\\[4pt]
The code said \bad{$\mathbf{+0.956 \pm 0.274}$}.
\end{center}

\bigskip
Three defects, all in arm 3's sampler. The decisive one: new points were
initialised \emph{off} the constraint surface $\{g(h)=y\}$, which the sampler can
never travel to once the likelihood is tight. Residual grew over iterations:

\begin{center}
\small
$0.0001 \to 0.004 \to 0.027 \to 0.019 \to 0.043 \to \mathbf{0.121}$
\end{center}

\bigskip
After the fix: \good{$-0.096 \pm 0.175$}, $\Ulost = 0$ exactly.

\end{frame}

% ------------------------------------------------------------------
\begin{frame}{Results: three predictions held}

\begin{center}
\begin{tabular}{llr}
\toprule
Prediction & Condition & Measured \\
\midrule
observing $h$ $=0$ exactly & $z$, $e^z$, $\sinh z$ & \good{$-0.000\pm0.000$} \\
knowing $g$ $\approx0$ & $y_1+y_2$ (linear) & \good{$+0.117\pm0.340$} \\
censoring away from optimum $\approx0$ & $-\max(z,0)$ & \good{$-0.066\pm0.209$} \\
censoring at the optimum $>0$ & clipped pair & \good{$+1.108\pm0.421$} \\
\bottomrule
\end{tabular}
\end{center}

\bigskip
The linear row is the one that bears on the paper: \textbf{on Truss and RCM40
there is no knowledge to exploit}, so whatever composite gains there must come
from the extra observations alone.

\end{frame}

% ------------------------------------------------------------------
\begin{frame}{The finding: \emph{where}, not \emph{how much}}

Both maps censor. \texttt{relu} destroys \textbf{more} --- it collapses an entire
half-line. Yet:

\begin{center}
\begin{tabular}{llrr}
\toprule
Map & Censoring sits & $\Ulost$ & observing $h$ \\
\midrule
$-\max(z,0)$ & away from the optimum & $0.075$ & \good{$-0.066\pm0.209$} \\
clipped pair & \textbf{at} the optimum & $0.580$ & \good{$+1.108\pm0.421$} \\
\bottomrule
\end{tabular}
\end{center}

\bigskip
The rest of the table agrees --- sorting by $\Ulost$ does not sort by gain:

\begin{center}
\small
\texttt{product} destroys $1.64$, pays $+0.21$ \quad$\cdot$\quad
\texttt{sum} destroys $0.70$, pays $-0.18$ \quad$\cdot$\quad
\texttt{abs} destroys $0.47$, pays $-0.06$
\end{center}

\bigskip
\begin{block}{Sharper hypothesis}
The question for a new problem is not ``is $g$ invertible?'' but
\textbf{``is $g$ invertible near the Pareto front?''}
\end{block}

\end{frame}

% ------------------------------------------------------------------
\begin{frame}{What did not hold: the sweep}

$(z-c)^2$ --- same form, one parameter, information loss tuned by $c$. Designed
to be the cleanest experiment in the set.

\bigskip
\begin{center}
\begin{tabular}{lrrrrr}
\toprule
$c$ & $0$ & $0.5$ & $1$ & $2$ & $3$ \\
\midrule
$\Ulost$ \textbf{at iteration 0} & $1.026$ & --- & $0.869$ & $0.104$ & $0.013$ \\
observing $h$ (predicted) & large & & $\downarrow$ & $\downarrow$ & $\to 0$ \\
observing $h$ \textbf{(measured)} & \bad{$+1.49$} & \bad{$+1.28$} & \bad{$+1.89$} & \bad{$+2.75$} & \bad{$+2.51$} \\
\bottomrule
\end{tabular}
\end{center}

\bigskip
The knob \emph{does} turn on the initial design. Then, over BO iterations,
$\Ulost$ falls for small $c$ and \textbf{rises} for large $c$ --- more data
cannot increase uncertainty about points already observed.

\medskip
Second symptom: ``knowing $g$'' goes \bad{negative} ($-1.18$ at $c=3$). Arm 3
holds the true kernel \emph{and} knows $g$; it should not lose to a plain GP.

\end{frame}

% ------------------------------------------------------------------
\begin{frame}{Diagnosis so far}

\begin{center}
\begin{tabular}{lll}
\toprule
Hypothesis & Test & Outcome \\
\midrule
Sampler budget too small & $4$ chains/$1500$ kept vs.\ $2$/$200$ & \good{ruled out} \\
Warm starts accumulate error & cold vs.\ warm, same design & \good{ruled out} \\
Level-set moves fail & acceptance in loop vs.\ in checks & \alert{confirmed} \\
\bottomrule
\end{tabular}
\end{center}

\bigskip
Move acceptance: $0.5$--$0.68$ on spread designs, \textbf{$0.000$} on clustered
ones. Flipping one point's sign against its packed neighbours is the
$10^{16}$-to-one rejection from the two-point example. The sampler stops changing its mind
exactly when BO tightens the design.

\medskip
\emph{Fix: block moves --- flip a whole neighbourhood at once.}

\bigskip
\begin{alertblock}{But that is not the whole story}
\texttt{identity} with the sampler forced on gives ``observing $h$''
$=+0.409\pm0.133$ with $\Ulost=0.0016$ --- it \emph{has} the information and
still loses. Raising the MC budget $16\times$ only moves it to $+0.275$.
\end{alertblock}

\end{frame}

% ------------------------------------------------------------------
\begin{frame}{Analogs vs.\ the benchmark results}

\begin{center}
\begin{tabular}{llrl}
\toprule
Analog & Benchmark & observing $h$ & Composite gain in the paper \\
\midrule
\texttt{product} & DTLZ2        & $+0.213$ & largest ($121\%$) \\
\texttt{max}     & Lake         & $+0.139$ & moderate \\
\texttt{ratio}   & RGB          & $+0.117$ & none \\
\texttt{norm}    & Color        & $-0.095$ & moderate \\
\texttt{sum}     & Truss, RCM40 & $-0.180$ & weakest \\
\bottomrule
\end{tabular}
\end{center}

\bigskip
Ordering matches at both ends.

\medskip
\textbf{Caveat, stated plainly:} standard errors are $0.06$--$0.25$, so no
individual pair separates significantly, and the correction subtracted from all
of them is itself uncertain. Consistent with the benchmarks --- not confirmation
of them.

\end{frame}

% ------------------------------------------------------------------
\begin{frame}{Where this goes next}

\textbf{To fix, in order}
\begin{enumerate}
\item The $+0.275$ residual on \texttt{identity} --- smallest case that still
shows the defect: one intermediate, no multimodality, near-point posterior.
Everything else is downstream.
\item Block level-set moves.
\item All-MLE hyperparameters in every arm (removes the last asymmetry rather
than correcting for it).
\end{enumerate}

\bigskip
\textbf{Cheapest test of the finding} --- no new BO runs
\begin{itemize}
\item Compute $\Ulost$ on each \emph{real} benchmark restricted to
near-Pareto-optimal points.
\item It should predict the composite gain better than $\Ulost$ over the whole
domain.
\item Uses points every existing run already evaluated.
\end{itemize}

\end{frame}

% ------------------------------------------------------------------
\begin{frame}{Summary}

\begin{itemize}
\item Composite BO bundles \textbf{knowing $g$} with \textbf{observing $h$};
a third arm separates them, and the three numbers must telescope.
\item Arm 3 is the hard part: no closed form, so MCMC over quantities never
observed.
\item \good{Held}: observing $h$ is exactly 0 when $g$ is injective; knowing $g$
is $\approx 0$ when $g$ is linear (Truss, RCM40).
\item \good{Finding}: \textbf{where} $g$ destroys information matters more than
how much --- censoring away from the optimum costs nothing, censoring at it
costs $13\times$ in final regret.
\item \bad{Open}: the sweep runs backwards; a $+0.275$ bias survives on a
control where the answer is provably 0.
\end{itemize}

\vfill
\centering
\small Code, results and full write-up:
\texttt{github.com/MuhammadAli-CS/three-arm-composite}

\end{frame}

\end{document}
